\chapter{L'entrepôt : architecture et garanties}
\label{ch:entrepot}

\questionchapitre{Comment transformer quatre jeux de données hétérogènes, volumineux et
piégeux en une ressource interrogeable, reproductible et digne de confiance ?}

Ce chapitre porte le résultat principal du travail. Il ne s'agit pas d'un choix d'outils mais
d'un ensemble de propriétés à garantir, dont dépend la possibilité même de comparer deux
expériences.

\section{Le principe : trois couches de qualité croissante}

La donnée traverse trois couches, matérialisées par trois schémas d'une même base
PostgreSQL~\cite{postgresql}. Ce patron répond à une exigence : ne jamais perdre la donnée
telle qu'elle a été reçue, et rendre chaque transformation ultérieure explicite, versionnée et
rejouable.

\begin{figure}[H]
\centering
\begin{tikzpicture}[
  font=\sffamily,
  src/.style={draw=black!55, fill=white, minimum width=2.45cm, minimum height=0.95cm,
              align=center, font=\sffamily\scriptsize},
  couche/.style={draw=black!80, fill=white, minimum width=9.6cm, minimum height=1.25cm,
                 align=left, font=\sffamily\small, inner xsep=10pt},
  cote/.style={draw=black!55, fill=black!4, minimum width=4.1cm, minimum height=1.25cm,
               align=left, font=\sffamily\scriptsize, inner xsep=8pt},
  fl/.style={-{Stealth[length=2.6mm]}, draw=black!75, very thick},
  outil/.style={font=\sffamily\scriptsize\itshape, color=black!65},
]

% --- Sources -------------------------------------------------------------
\node[src] (s1) {Hub'Eau\\piézométrie};
\node[src, right=0.28cm of s1] (s2) {Hub'Eau\\hydrométrie};
\node[src, right=0.28cm of s2] (s3) {Copernicus\\ERA5-Land};
\node[src, right=0.28cm of s3] (s4) {BDLISA};

% --- Couches -------------------------------------------------------------
\node[couche] (bronze) at ($(s1.south)!0.5!(s4.south)+(0,-2.05)$)
  {\textbf{Bronze} \quad \scriptsize charges utiles brutes, colonnes texte\\[2pt]
   \scriptsize\itshape on n'interprète rien};

\node[couche, below=1.15cm of bronze] (silver)
  {\textbf{Silver} \quad \scriptsize typé, nettoyé, dédupliqué\\[2pt]
   \scriptsize\itshape toute ligne écartée garde le motif de son exclusion};

\node[couche, below=1.15cm of silver] (gold)
  {\textbf{Gold} \quad \scriptsize stations jointes au climat et au référentiel\\[2pt]
   \scriptsize\itshape interrogeable en SQL, sans rien savoir de l'amont};

% --- Rejets, à côté de Silver -------------------------------------------
\node[cote, right=0.5cm of silver] (rejets)
  {\textbf{Rejets}\\[2pt] date absente, code nul,\\station orpheline\ldots};

% --- Consommateurs -------------------------------------------------------
\node[src, below=1.1cm of gold.south, anchor=north, minimum width=5.2cm, minimum height=0.8cm] (conso)
  {Consommateurs : plateforme,\\carnets de calcul, outils tiers};

% --- Flèches -------------------------------------------------------------
% Les quatre sources descendent sur une ligne de collecte commune, d'où part
% une flèche unique : le trajet reste lisible quel que soit le nombre de sources.
\coordinate (collecte) at ($(s1.south)!0.5!(s4.south)+(0,-0.5)$);
\foreach \s in {s1,s2,s3,s4} \draw[draw=black!75, very thick] (\s.south) -- (\s.south |- collecte);
\draw[draw=black!75, very thick] (s1.south |- collecte) -- (s4.south |- collecte);
\draw[fl] (collecte) -- node[outil, right=3pt]{dlt} (bronze.north);

\draw[fl] (bronze) -- node[outil, right=3pt]{dbt} (silver);
\draw[fl] (silver) -- node[outil, right=3pt]{dbt} (gold);
\draw[fl] (gold) -- (conso);
\draw[fl, densely dashed, thick] (silver.east) -- (rejets.west);

\end{tikzpicture}
\caption{Les trois couches de l'entrepôt. La donnée brute n'est jamais écrasée, et les lignes
écartées au nettoyage sont conservées à part avec le motif de leur exclusion : le filtrage se
vérifie par une requête au lieu d'être silencieux.}
\label{fig:medaillon}
\end{figure}

\section{Ingestion : ne rien décider}

Chaque source est ingérée par une bibliothèque d'extraction déclarative~\cite{dlt} qui prend en
charge la pagination, les tentatives successives avec temporisation exponentielle sur erreur
serveur, l'inférence de schéma et la déduplication par fusion sur clé primaire. Les séries
temporelles sont partitionnées par année, chaque tranche restant rechargeable indépendamment.
Les partitions couvrent la période 1967 à aujourd'hui pour la piézométrie comme pour
l'hydrométrie, la profondeur effectivement disponible variant d'une station à l'autre.

Le principe de la couche est de ne rien interpréter : les colonnes restent textuelles, les
valeurs sont celles reçues. Toute interprétation appartient aux couches suivantes, où elle est
lisible et modifiable.

Le cas d'ERA5 fait exception sur un point imposé par la nature du produit. Les fichiers NetCDF
ne sont jamais archivés en base : ils sont téléchargés en stockage temporaire, extraits en
mémoire~\cite{xarray} et insérés directement. Deux chemins d'ingestion coexistent, l'un pour
les séries issues du pas de \SI{00}{\hour}~UTC, l'autre pour les statistiques journalières de
température agrégées localement sur les vingt-quatre pas horaires, pour la raison exposée au
chapitre~\ref{ch:verrou}.

\begin{center}
\captionof{table}{Volumétrie de la couche brute, mesurée le 24 août 2026.}
\begin{tabularx}{\textwidth}{@{}Xr@{}}
\toprule
\textbf{Table} & \textbf{Lignes} \\
\midrule
Statistiques journalières de température ERA5 & $\approx$ \num{320000000} \\
Séries temporelles ERA5 & $\approx$ \num{300000000} \\
Mesures piézométriques & $\approx$ \num{23000000} \\
Observations hydrométriques élaborées & $\approx$ \num{15000000} \\
Ouvrages piézométriques & \num{23333} \\
Sites hydrométriques & \num{9283} \\
Stations hydrométriques & \num{6468} \\
Entités hydrogéologiques BDLISA & \num{3716} \\
\midrule
\textbf{Total} & $\approx$ \num{658000000} \\
\bottomrule
\end{tabularx}
\end{center}

\section{Nettoyage : filtrer en rendant compte}

Chaque modèle de la couche intermédiaire~\cite{dbt} sélectionne les colonnes utiles, convertit
les types, déduplique et écarte les lignes invalides. La caractéristique importante n'est pas
le nettoyage mais sa traçabilité : les lignes écartées ne disparaissent pas, elles sont écrites
dans un schéma dédié avec le motif de leur exclusion, par exemple une date de mesure absente,
un code de station manquant ou une station orpheline de son site.

Cette décision a un coût de stockage négligeable et une valeur méthodologique élevée. Un
filtrage silencieux est invérifiable : rien ne distingue une station réellement absente d'une
station éliminée par un critère trop strict. Ici, la question se tranche par une requête.

Huit modèles de nettoyage et trois modèles de rejets composent cette couche. Les modèles de
séries temporelles sont incrémentaux, avec une fenêtre de recouvrement de sept jours qui
absorbe les corrections tardives publiées par les producteurs.

\section{Couche analytique : joindre, enrichir, exposer}

C'est ici que se résout le verrou central du chapitre précédent.

\paragraph{Le rattachement climatique.} Chaque station est associée à sa maille ERA5 la plus
proche par une jointure spatiale du plus proche voisin, exécutée par l'extension
géospatiale~\cite{postgis} au moyen d'une jointure latérale et de l'opérateur de distance
indexé. Cette opération, effectuée une fois et matérialisée, transforme deux jeux de données
étrangers l'un à l'autre en une chronique unique où le niveau et son forçage climatique
figurent sur la même ligne, au même jour.

\paragraph{L'enrichissement hydrogéologique.} Les entités du référentiel BDLISA sont rattachées
aux stations, rendant disponible la nature de l'aquifère observé, information que le
chapitre~\ref{ch:metier} exploitera pour paramétrer les modèles métier.

\paragraph{Les tables finales.} La couche expose des faits journaliers, des faits agrégés
mensuels et annuels, des dimensions de description (calendrier, géographie, stations) et les
grilles climatiques, soit dix-huit modèles de transformation complétés par des tables
d'indicateurs calculées en Python. Les faits journaliers sont stockés en tables partitionnées
dans le temps et compressées~\cite{timescaledb}, ce qui maintient les performances
d'interrogation sur plusieurs dizaines de millions de lignes.

L'ensemble est interrogeable en SQL. C'est le contrat d'interface : un consommateur n'a besoin
de connaître ni la chaîne d'ingestion, ni les outils employés, seulement le schéma des tables
finales.

\section{Orchestration : planifier l'ingestion, déclencher la transformation}

L'ingestion est \textbf{planifiée}~\cite{dagster} : mise à jour quotidienne des chroniques et
du climat, rafraîchissement mensuel du référentiel, contrôle hebdomadaire de complétude,
recalcul hebdomadaire des références statistiques. La transformation est \textbf{déclenchée par
événement} : dès que de nouvelles données brutes sont disponibles, un capteur lance la chaîne
de transformation, laquelle déclenche à son tour le recalcul des indicateurs et les contrôles
de qualité. L'ensemble compte huit traitements planifiés et trois capteurs.

Ce choix évite le défaut des chaînes entièrement planifiées, où la transformation s'exécute à
heure fixe sur des données parfois absentes et produit des résultats partiels sans que rien ne
le signale.

\begin{remarque}[title={Les contrôles de qualité ne bloquent pas le rafraîchissement}]
Les tests s'exécutent en parallèle du recalcul des indicateurs, et non en amont. Un test en
échec fait échouer sa propre exécution, donc alerte, mais n'empêche pas la donnée d'être mise à
jour. L'arbitrage est délibéré : sur une chaîne de surveillance, une donnée imparfaite mais
fraîche vaut mieux qu'une donnée absente, à condition que le défaut soit signalé.
\end{remarque}

\section{Les garanties recherchées, et ce que leur mise à l'épreuve a coûté}

Trois propriétés ont été recherchées, puis vérifiées par l'exploitation réelle. Les incidents
qui suivent sont rapportés parce qu'ils constituent le contenu concret du travail d'ingénierie :
une architecture n'est pas correcte parce qu'elle est bien dessinée, elle le devient une fois
ses modes de défaillance rencontrés et traités.

\subsection{Idempotence}

Relancer une ingestion sur une période déjà chargée ne doit rien dupliquer. La propriété est
obtenue par fusion sur clé primaire à l'ingestion, et par suppression puis réinsertion de la
plage concernée pour les chargements historiques.

Elle a été mise en défaut de deux manières. D'abord par six modèles dont le prédicat de
suppression incrémentale était ancré sur la date courante alors que leur sélection l'était sur
le maximum de la table : sur un jeu plus ancien que la fenêtre, la suppression ne trouvait rien
et l'insertion répétait des lignes déjà présentes, jusqu'à violation de contrainte d'unicité.
Le prédicat, qui n'était qu'une optimisation d'élagage, a été retiré plutôt que réparé.
Ensuite par le lancement d'une matérialisation depuis la ligne de commande, qui contourne la
file d'exécution : le client interrompu, le traitement continuait d'insérer côté serveur. Des
doublons ont été observés ainsi le 7 juillet 2026.

\subsection{Maîtrise de la concurrence}

Les exécutions ne sont pas globalement sérialisées, ce qui serait simple mais inutilement lent.
Elles sont contraintes par clés de concurrence, chaque clé plafonnant le nombre d'exécutions
simultanées d'une même nature : une seule ingestion par source à la fois, une seule chaîne de
transformation, mais trois partitions climatiques en vol pour le débit.

Une de ces clés est partagée entre la mise à jour nocturne des températures et le rechargement
historique. La table concernée ne porte pas de contrainte d'unicité et son remplacement n'est
pas atomique : lorsque la partition de l'année courante recouvrait la fenêtre nocturne, les
deux traitements dupliquaient des lignes. Le partage de clé rend ce recouvrement impossible, au
prix d'un seul créneau d'exécution sur cinq, de sorte que le rechargement historique n'affame
jamais la production.

\subsection{Reprise}

Le chargement initial complet est interruptible et redémarrable, son avancement étant persisté
en base. C'est une nécessité pratique : le chargement se compte en jours et en dizaines de
gigaoctets, et aucune exécution de cette durée ne se déroule sans incident.

\subsection{Les défaillances silencieuses, et pourquoi elles sont les plus coûteuses}

Trois situations produisent des résultats faux sans émettre la moindre erreur. Elles sont
documentées dans le dépôt parce qu'elles sont invisibles à l'exécution et coûteuses à
diagnostiquer.

\begin{itemize}
  \item \textbf{Une colonne entièrement nulle.} Si les séries climatiques sont chargées avant
  les statistiques journalières de température pour la même période, la garde incrémentale du
  modèle intermédiaire ne trouve aucune ligne nouvelle, la jointure externe rend des valeurs
  nulles, et la température est absente de toutes les lignes. Aucun test n'échoue, puisque la
  jointure est externe et qu'aucune contrainte ne porte sur la colonne. Mesuré le 24 août
  2026 : \SI{0}{\percent} de lignes renseignées en température, contre \SI{95,4}{\percent}
  après reconstruction complète du modèle intermédiaire.
  \item \textbf{Un rechargement historique qui n'atteint pas la couche intermédiaire.} Le
  modèle de nettoyage des températures fonctionne par ajout avec un filtre sur son propre
  maximum. Le traitement nocturne ayant déjà renseigné l'année courante, toute ligne
  rétrospective est silencieusement ignorée.
  \item \textbf{Un modèle incrémental qui ne produit rien.} Sur un jeu de données arrêté dans
  le passé, la fenêtre de recalcul glissante ne couvre aucune donnée. L'exécution est déclarée
  réussie et la table ne grandit pas.
\end{itemize}

\begin{remarque}[title={Ce que ces incidents enseignent}]
Une chaîne correcte en régime nominal peut être fausse en régime de recouvrement ou de
rattrapage. Les garanties d'unicité doivent être portées par le schéma quand c'est possible et
par l'orchestration quand ce ne l'est pas, jamais par l'hypothèse implicite que deux
traitements ne se croiseront pas. Surtout, une vérification qui ne peut pas échouer ne vérifie
rien : les trois cas ci-dessus passent tous les tests existants.
\end{remarque}

\section{Infrastructure}

Sept services conteneurisés composent le déploiement. Le code métier et l'orchestrateur sont
séparés en deux images communiquant par appel de procédure distante : l'image métier, lourde,
embarque le code Python et les bibliothèques géospatiales ; l'image d'orchestration, légère, ne
contient aucun code métier et n'est reconstruite que lors d'un changement de version de
l'ordonnanceur. Les volumes de données sont déclarés hors du cycle de vie des conteneurs, de
sorte qu'une commande d'arrêt ne peut pas les détruire.

\begin{acquisouvert}
\textbf{Acquis.} Une chaîne complète, automatisée, idempotente, documentée et éprouvée par
l'exploitation, dont les tables finales sont consommables en SQL sans connaissance de l'amont.
Le coût d'entrée à l'expérimentation passe de plusieurs semaines à une requête.

\textbf{Ouvert.} L'entrepôt couvre la quantité d'eau ; l'extension à la qualité supposerait de
nouvelles sources et de nouveaux modèles, sans remettre en cause l'architecture. La supervision
de production reste à instrumenter au-delà des contrôles existants.
\end{acquisouvert}
